\chapter{Correlation of Projection Differences with Behaviour For Each Step}
\newpage

\begin{table}[htbp]
    \centering
    \scriptsize
    \setlength{\tabcolsep}{3pt}
    % Generated by method.visualization.make_plots -- do not edit.
\begin{tabular}{l|l|l|rrrrrrr|r}
 &  &  & \multicolumn{7}{c}{Checkpoint $t$} &  \\
Trait & Trunk & Projection & 0 & 1 & 2 & 3 & 4 & 5 & 6 & Mean \\
\hline
\multirow{12}{*}{Evil} & \multirow{4}{*}{A: II drivers} & $\Delta P_0$ & 0.99 & \textbf{0.96} & \textbf{0.99} & \textbf{0.86} & 0.86 & 0.85 & 0.85 & 0.91 \\
 &  & $\Delta \hat{P}_t^{(\mathbf{v}_0)}$ & 0.99 & 0.95 & 0.98 & 0.80 & 0.79 & 0.77 & 0.76 & 0.86 \\
 &  & $\Delta \hat{P}_t$ & 0.99 & 0.93 & 0.98 & 0.78 & 0.78 & 0.76 & 0.74 & 0.85 \\
 &  & $\Delta P_t$ & 0.99 & 0.94 & 0.96 & 0.83 & \textbf{0.89} & \textbf{0.90} & \textbf{0.93} & \textbf{0.92} \\
\cline{2-11}
 & \multirow{4}{*}{B: I drivers} & $\Delta P_0$ & 0.99 & \textbf{0.89} & 0.92 & \textbf{0.95} & \textbf{0.96} & \textbf{0.95} & \textbf{0.99} & \textbf{0.95} \\
 &  & $\Delta \hat{P}_t^{(\mathbf{v}_0)}$ & 0.99 & 0.84 & 0.86 & 0.92 & 0.91 & 0.88 & \textbf{0.99} & 0.91 \\
 &  & $\Delta \hat{P}_t$ & 0.99 & 0.82 & 0.85 & 0.91 & 0.90 & 0.88 & \textbf{0.99} & 0.91 \\
 &  & $\Delta P_t$ & 0.99 & 0.81 & \textbf{0.95} & 0.90 & 0.88 & 0.87 & 0.97 & 0.91 \\
\cline{2-11}
 & \multirow{4}{*}{C: Normal only (control)} & $\Delta P_0$ & 0.99 & \textbf{0.99} & \textbf{0.99} & \textbf{0.98} & \textbf{0.99} & \textbf{0.99} & \textbf{0.99} & \textbf{0.99} \\
 &  & $\Delta \hat{P}_t^{(\mathbf{v}_0)}$ & 0.99 & \textbf{0.99} & \textbf{0.99} & \textbf{0.98} & \textbf{0.99} & \textbf{0.99} & \textbf{0.99} & \textbf{0.99} \\
 &  & $\Delta \hat{P}_t$ & 0.99 & \textbf{0.99} & 0.98 & \textbf{0.98} & \textbf{0.99} & \textbf{0.99} & \textbf{0.99} & \textbf{0.99} \\
 &  & $\Delta P_t$ & 0.99 & 0.97 & 0.97 & 0.97 & 0.98 & 0.98 & 0.96 & 0.98 \\
\hline
\multirow{12}{*}{Sycophancy} & \multirow{4}{*}{A: II drivers} & $\Delta P_0$ & 0.86 & 0.51 & 0.85 & 0.63 & 0.72 & 0.23 & \textbf{0.29} & 0.58 \\
 &  & $\Delta \hat{P}_t^{(\mathbf{v}_0)}$ & 0.86 & 0.53 & 0.90 & 0.70 & 0.77 & 0.15 & 0.25 & 0.59 \\
 &  & $\Delta \hat{P}_t$ & 0.86 & 0.48 & 0.89 & 0.62 & 0.73 & 0.10 & 0.20 & 0.55 \\
 &  & $\Delta P_t$ & 0.86 & \textbf{0.73} & \textbf{0.92} & \textbf{0.88} & \textbf{0.78} & \textbf{0.79} & 0.27 & \textbf{0.75} \\
\cline{2-11}
 & \multirow{4}{*}{B: I drivers} & $\Delta P_0$ & 0.86 & 0.82 & 0.72 & 0.88 & 0.82 & 0.78 & 0.92 & 0.83 \\
 &  & $\Delta \hat{P}_t^{(\mathbf{v}_0)}$ & 0.86 & 0.89 & 0.82 & \textbf{0.93} & 0.89 & 0.85 & 0.95 & 0.88 \\
 &  & $\Delta \hat{P}_t$ & 0.86 & 0.87 & 0.81 & 0.92 & 0.87 & 0.84 & 0.95 & 0.87 \\
 &  & $\Delta P_t$ & 0.86 & \textbf{0.94} & \textbf{0.83} & 0.92 & \textbf{0.95} & \textbf{0.93} & \textbf{0.96} & \textbf{0.91} \\
\cline{2-11}
 & \multirow{4}{*}{C: Normal only (control)} & $\Delta P_0$ & 0.86 & 0.90 & 0.89 & 0.91 & 0.88 & 0.92 & 0.86 & 0.89 \\
 &  & $\Delta \hat{P}_t^{(\mathbf{v}_0)}$ & 0.86 & \textbf{0.93} & \textbf{0.92} & \textbf{0.95} & \textbf{0.92} & 0.95 & 0.88 & \textbf{0.92} \\
 &  & $\Delta \hat{P}_t$ & 0.86 & \textbf{0.93} & \textbf{0.92} & \textbf{0.95} & \textbf{0.92} & 0.95 & 0.87 & 0.91 \\
 &  & $\Delta P_t$ & 0.86 & \textbf{0.93} & 0.91 & 0.92 & 0.90 & \textbf{0.96} & \textbf{0.94} & \textbf{0.92} \\
\end{tabular}

    \caption{Correlation $r$ between each projection difference and the behaviour reached after the step, at every checkpoint of every trunk for each of the projection difference variants discussed in \Cref{sec:deltap-variants}.}
    \label{tab:decay-correlations}
\end{table}

% \begin{table}[H]
%     \centering
%     \scriptsize
%     \setlength{\tabcolsep}{3pt}
%     % Generated by method.visualization.make_plots -- do not edit.
% BIAS in judge points; forecast: $M_0$ fit on $\Delta b$.
\begin{tabular}{l|l|l|rrrrrrr|r}
 &  &  & \multicolumn{7}{c}{Checkpoint $t$} &  \\
Trait & Trunk & Projection & 0 & 1 & 2 & 3 & 4 & 5 & 6 & Mean \\
\hline
\multirow{21}{*}{Evil} & \multirow{7}{*}{A: II drivers} & $\Delta P_0$ & -2.2 & -2.9 & 0.3 & 36.1 & 3.6 & 20.4 & 14.3 & 9.9 \\
 &  & $\Delta P_t^{0\leftarrow 0,[0]}$ & -2.2 & -1.5 & \textbf{-0.2} & 38.0 & 4.3 & 20.0 & 13.9 & 10.3 \\
 &  & $\Delta P_t^{t\leftarrow 0,[0]}$ & -2.2 & \textbf{-0.1} & 0.6 & 41.1 & 6.6 & 23.0 & 16.8 & 12.3 \\
 &  & $\Delta P_t^{t\leftarrow t,[0]}$ & \textbf{-2.1} & 0.7 & -2.7 & 34.2 & \textbf{-1.8} & 15.7 & 9.9 & 7.7 \\
 &  & $\Delta P_t^{0\leftarrow 0,[t]}$ & -2.2 & -22.5 & -4.6 & 6.6 & -12.0 & \textbf{-5.9} & \textbf{-5.8} & \textbf{-6.6} \\
 &  & $\Delta P_t^{t\leftarrow 0,[t]}$ & -2.2 & -23.7 & -4.2 & \textbf{4.0} & -12.2 & -7.8 & -6.4 & -7.5 \\
 &  & $\Delta P_t^{t\leftarrow t,[t]}$ & \textbf{-2.1} & -24.2 & -3.1 & 4.2 & -14.3 & -6.2 & -7.1 & -7.5 \\
\cline{2-11}
 & \multirow{7}{*}{B: I drivers} & $\Delta P_0$ & -2.2 & 23.2 & 19.0 & -3.9 & 10.9 & 9.2 & \textbf{-2.2} & 7.7 \\
 &  & $\Delta P_t^{0\leftarrow 0,[0]}$ & -2.2 & 27.4 & 19.7 & \textbf{-3.6} & 11.4 & 8.7 & -4.2 & 8.2 \\
 &  & $\Delta P_t^{t\leftarrow 0,[0]}$ & -2.2 & 26.8 & 19.7 & \textbf{-3.6} & 12.6 & 10.6 & -3.0 & 8.7 \\
 &  & $\Delta P_t^{t\leftarrow t,[0]}$ & \textbf{-2.1} & 12.2 & 11.3 & -12.5 & \textbf{2.6} & \textbf{0.7} & -10.4 & \textbf{0.3} \\
 &  & $\Delta P_t^{0\leftarrow 0,[t]}$ & -2.2 & 3.3 & 3.9 & -9.5 & -2.9 & -5.7 & -7.6 & -2.9 \\
 &  & $\Delta P_t^{t\leftarrow 0,[t]}$ & -2.2 & \textbf{2.7} & 3.5 & -9.1 & -2.9 & -5.7 & -7.2 & -3.0 \\
 &  & $\Delta P_t^{t\leftarrow t,[t]}$ & \textbf{-2.1} & 2.9 & \textbf{0.1} & -10.5 & -3.6 & -6.0 & -7.4 & -3.8 \\
\cline{2-11}
 & \multirow{7}{*}{C: Normal only} & $\Delta P_0$ & -2.2 & \textbf{0.2} & \textbf{-0.7} & \textbf{0.3} & \textbf{-0.7} & \textbf{-0.2} & \textbf{0.2} & \textbf{-0.4} \\
 &  & $\Delta P_t^{0\leftarrow 0,[0]}$ & -2.2 & -1.3 & -2.1 & -1.2 & -2.7 & -3.2 & -3.1 & -2.2 \\
 &  & $\Delta P_t^{t\leftarrow 0,[0]}$ & -2.2 & -1.9 & -2.5 & -1.6 & -2.8 & -2.6 & -1.8 & -2.2 \\
 &  & $\Delta P_t^{t\leftarrow t,[0]}$ & \textbf{-2.1} & -7.9 & -8.1 & -8.8 & -9.6 & -8.4 & -6.6 & -7.4 \\
 &  & $\Delta P_t^{0\leftarrow 0,[t]}$ & -2.2 & -3.5 & -5.9 & -6.3 & -7.4 & -5.5 & -5.0 & -5.1 \\
 &  & $\Delta P_t^{t\leftarrow 0,[t]}$ & -2.2 & -3.3 & -5.8 & -6.2 & -7.2 & -5.5 & -4.8 & -5.0 \\
 &  & $\Delta P_t^{t\leftarrow t,[t]}$ & \textbf{-2.1} & -2.4 & -4.5 & -4.7 & -6.2 & -4.4 & \textbf{-0.2} & -3.5 \\
\hline
\multirow{21}{*}{Sycophancy} & \multirow{7}{*}{A: II drivers} & $\Delta P_0$ & 3.5 & 43.6 & 5.5 & 35.6 & 18.5 & 53.2 & 49.6 & 29.9 \\
 &  & $\Delta P_t^{0\leftarrow 0,[0]}$ & 3.5 & 46.0 & 3.6 & 32.7 & 15.3 & 53.7 & 48.8 & 29.1 \\
 &  & $\Delta P_t^{t\leftarrow 0,[0]}$ & 3.5 & 51.2 & 5.1 & 45.0 & 24.3 & 63.9 & 59.1 & 36.0 \\
 &  & $\Delta P_t^{t\leftarrow t,[0]}$ & \textbf{3.0} & 48.2 & \textbf{-2.5} & 42.1 & 19.2 & 55.5 & 51.9 & 31.1 \\
 &  & $\Delta P_t^{0\leftarrow 0,[t]}$ & 3.5 & 12.7 & -7.1 & 7.8 & -0.7 & 9.7 & 18.2 & 6.3 \\
 &  & $\Delta P_t^{t\leftarrow 0,[t]}$ & 3.5 & 10.3 & -6.9 & \textbf{-0.6} & \textbf{0.3} & 2.9 & 15.3 & 3.5 \\
 &  & $\Delta P_t^{t\leftarrow t,[t]}$ & \textbf{3.0} & \textbf{10.2} & -4.2 & -7.0 & -3.9 & \textbf{1.9} & \textbf{11.7} & \textbf{1.7} \\
\cline{2-11}
 & \multirow{7}{*}{B: I drivers} & $\Delta P_0$ & 3.5 & 16.3 & 15.7 & \textbf{-1.2} & 10.5 & 18.3 & 3.7 & 9.5 \\
 &  & $\Delta P_t^{0\leftarrow 0,[0]}$ & 3.5 & 19.8 & 16.1 & -4.5 & 8.8 & 15.9 & -3.6 & 8.0 \\
 &  & $\Delta P_t^{t\leftarrow 0,[0]}$ & 3.5 & 21.4 & 18.2 & -4.5 & 11.9 & 19.9 & \textbf{-3.3} & 9.6 \\
 &  & $\Delta P_t^{t\leftarrow t,[0]}$ & \textbf{3.0} & \textbf{8.7} & 8.0 & -18.2 & \textbf{-1.4} & 5.8 & -17.4 & \textbf{-1.6} \\
 &  & $\Delta P_t^{0\leftarrow 0,[t]}$ & 3.5 & -9.7 & -4.4 & -9.1 & -6.5 & -1.3 & -4.3 & -4.5 \\
 &  & $\Delta P_t^{t\leftarrow 0,[t]}$ & 3.5 & -9.5 & \textbf{-4.0} & -8.3 & -6.0 & \textbf{-1.2} & -3.6 & -4.2 \\
 &  & $\Delta P_t^{t\leftarrow t,[t]}$ & \textbf{3.0} & -9.6 & -10.4 & -11.5 & -8.8 & -4.0 & -4.0 & -6.5 \\
\cline{2-11}
 & \multirow{7}{*}{C: Normal only} & $\Delta P_0$ & 3.5 & 5.2 & 3.3 & 3.4 & 5.1 & \textbf{3.6} & 4.5 & 4.1 \\
 &  & $\Delta P_t^{0\leftarrow 0,[0]}$ & 3.5 & 1.4 & \textbf{-2.2} & \textbf{-2.5} & \textbf{-1.8} & -5.0 & -3.3 & \textbf{-1.4} \\
 &  & $\Delta P_t^{t\leftarrow 0,[0]}$ & 3.5 & \textbf{-0.4} & -3.7 & -4.6 & -2.9 & -4.9 & \textbf{-1.5} & -2.1 \\
 &  & $\Delta P_t^{t\leftarrow t,[0]}$ & \textbf{3.0} & -10.6 & -15.6 & -18.3 & -16.7 & -15.9 & -10.0 & -12.0 \\
 &  & $\Delta P_t^{0\leftarrow 0,[t]}$ & 3.5 & -7.6 & -5.9 & -7.9 & -6.5 & -4.9 & -4.5 & -4.8 \\
 &  & $\Delta P_t^{t\leftarrow 0,[t]}$ & 3.5 & -8.4 & -6.7 & -8.8 & -6.8 & -5.0 & -4.2 & -5.2 \\
 &  & $\Delta P_t^{t\leftarrow t,[t]}$ & \textbf{3.0} & -7.3 & -5.8 & -7.3 & -6.5 & -3.9 & 6.2 & -3.1 \\
\end{tabular}

%     \caption{Mean signed error, in judge points, of the same frozen forecasts:
%     positive where $f_0$ promised more behaviour change than the step
%     delivered. This is what staleness looks like that an RMSE cannot separate
%     from scatter --- as a trunk drifts, the frozen line over-predicts every
%     probe at once. The refitted forecaster is omitted rather than shown at
%     zero: least squares with a free intercept leaves residuals summing to zero,
%     so its bias is an identity and not a measurement. Bold marks the row
%     nearest zero at each checkpoint, over- and under-prediction counting alike.}
%     \label{tab:forecast-bias}
% \end{table}


% \begin{table}[H]
%     \centering
%     \small
%     \setlength{\tabcolsep}{3pt}
%     % Generated by method.visualization.make_plots -- do not edit.
% BIAS in judge points; projection: $\Delta P_0$.
\begin{tabular}{l|l|l|rrrrrrr|r}
 &  &  & \multicolumn{7}{c}{Checkpoint $t$} &  \\
Trait & Trunk & Forecast & 0 & 1 & 2 & 3 & 4 & 5 & 6 & Mean \\
\hline
\multirow{21}{*}{Evil} & \multirow{7}{*}{A: II drivers} & $M_0$ fit on $\Delta b$ & -2.2 & -2.9 & 0.3 & 36.1 & 3.6 & 20.4 & 14.3 & 9.9 \\
 &  & $f_0$ $\times\, c_t(\mathbf{z}_t^{[0,0]})$ & \textbf{-0.1} & -24.9 & -2.2 & 12.5 & -3.3 & \textbf{-1.6} & \textbf{1.2} & -2.6 \\
 &  & $f_0$ $\times\, c_t(p_t^{[0]})$ & 1.3 & -16.1 & \textbf{0.2} & 18.7 & -1.5 & 2.6 & 3.0 & \textbf{1.2} \\
 &  & $f_0$ $\times\, c_t(q_t^{[0,0]})$ & 1.6 & -9.2 & -3.2 & 28.9 & -1.7 & 13.7 & 7.7 & 5.4 \\
 &  & $f_0$ $\times\, c_t(\rho_t^{[0]})$ & -0.6 & -5.3 & -2.0 & 30.1 & -2.0 & 13.0 & 7.2 & 5.8 \\
 &  & $f_0$ $\times\, c_t(r_t^{[0]})$ & -0.3 & \textbf{-1.3} & -0.3 & 33.0 & \textbf{-0.3} & 15.1 & 8.5 & 7.8 \\
 &  & $f_0$ $\times\, c_t(b_t)$ & -1.5 & -3.8 & 1.0 & \textbf{12.1} & -3.4 & 5.3 & 2.2 & 1.7 \\
\cline{2-11}
 & \multirow{7}{*}{B: I drivers} & $M_0$ fit on $\Delta b$ & -2.2 & 23.2 & 19.0 & \textbf{-3.9} & 10.9 & 9.2 & -2.2 & 7.7 \\
 &  & $f_0$ $\times\, c_t(\mathbf{z}_t^{[0,0]})$ & \textbf{-0.1} & \textbf{-1.9} & -1.9 & -14.2 & -3.0 & -7.5 & \textbf{-1.6} & -4.3 \\
 &  & $f_0$ $\times\, c_t(p_t^{[0]})$ & 1.0 & 5.8 & 3.0 & -7.8 & 2.7 & \textbf{-1.3} & -2.1 & \textbf{0.2} \\
 &  & $f_0$ $\times\, c_t(q_t^{[0,0]})$ & 4.5 & 15.5 & 5.9 & -12.8 & -1.5 & -4.0 & -12.4 & -0.7 \\
 &  & $f_0$ $\times\, c_t(\rho_t^{[0]})$ & 3.9 & 17.6 & 9.3 & -11.7 & -1.1 & -4.6 & -14.7 & \textbf{-0.2} \\
 &  & $f_0$ $\times\, c_t(r_t^{[0]})$ & 0.8 & 13.7 & 7.0 & -18.0 & -5.1 & -7.9 & -14.4 & -3.4 \\
 &  & $f_0$ $\times\, c_t(b_t)$ & -1.6 & \textbf{1.9} & \textbf{0.5} & -6.9 & \textbf{-1.0} & -1.4 & \textbf{-1.6} & -1.5 \\
\cline{2-11}
 & \multirow{7}{*}{C: Normal only} & $M_0$ fit on $\Delta b$ & -2.2 & \textbf{0.2} & -0.7 & 0.3 & \textbf{-0.7} & \textbf{-0.2} & \textbf{0.2} & -0.4 \\
 &  & $f_0$ $\times\, c_t(\mathbf{z}_t^{[0,0]})$ & -0.6 & 3.5 & 2.5 & 1.3 & -0.8 & 2.1 & -7.8 & \textbf{0.0} \\
 &  & $f_0$ $\times\, c_t(p_t^{[0]})$ & 2.2 & 1.7 & \textbf{0.0} & \textbf{0.2} & -1.5 & 1.7 & -1.5 & 0.4 \\
 &  & $f_0$ $\times\, c_t(q_t^{[0,0]})$ & 0.7 & -2.9 & -5.8 & -4.9 & -6.6 & -6.0 & -4.3 & -4.3 \\
 &  & $f_0$ $\times\, c_t(\rho_t^{[0]})$ & \textbf{-0.2} & -1.9 & -4.6 & -4.8 & -6.7 & -7.1 & -8.9 & -4.9 \\
 &  & $f_0$ $\times\, c_t(r_t^{[0]})$ & -4.2 & -4.0 & -7.1 & -5.7 & -7.6 & -6.9 & -8.5 & -6.3 \\
 &  & $f_0$ $\times\, c_t(b_t)$ & 0.4 & 2.7 & 1.9 & 2.9 & 1.8 & 2.3 & 2.8 & 2.1 \\
\hline
\multirow{21}{*}{Sycophancy} & \multirow{7}{*}{A: II drivers} & $M_0$ fit on $\Delta b$ & 3.5 & 43.6 & 5.5 & 35.6 & 18.5 & 53.2 & 49.6 & 29.9 \\
 &  & $f_0$ $\times\, c_t(\mathbf{z}_t^{[0,0]})$ & \textbf{-0.4} & 24.8 & -1.8 & 15.1 & 7.5 & 15.0 & 20.4 & 11.5 \\
 &  & $f_0$ $\times\, c_t(p_t^{[0]})$ & -4.7 & 27.1 & -2.8 & 26.2 & 15.0 & 33.8 & 38.7 & 19.0 \\
 &  & $f_0$ $\times\, c_t(q_t^{[0,0]})$ & 0.8 & 37.0 & \textbf{-0.3} & 28.2 & 11.7 & 45.1 & 41.6 & 23.5 \\
 &  & $f_0$ $\times\, c_t(\rho_t^{[0]})$ & 1.1 & 38.0 & 1.0 & 24.3 & 8.6 & 38.1 & 35.3 & 20.9 \\
 &  & $f_0$ $\times\, c_t(r_t^{[0]})$ & -3.8 & 35.5 & -0.5 & 29.2 & 13.9 & 43.9 & 41.3 & 22.8 \\
 &  & $f_0$ $\times\, c_t(b_t)$ & 0.7 & \textbf{14.8} & 1.4 & \textbf{11.1} & \textbf{2.8} & \textbf{10.4} & \textbf{9.6} & \textbf{7.2} \\
\cline{2-11}
 & \multirow{7}{*}{B: I drivers} & $M_0$ fit on $\Delta b$ & 3.5 & 16.3 & 15.7 & -1.2 & 10.5 & 18.3 & 3.7 & 9.5 \\
 &  & $f_0$ $\times\, c_t(\mathbf{z}_t^{[0,0]})$ & 1.5 & \textbf{-0.4} & -3.6 & -7.6 & -4.3 & -2.2 & -3.9 & -2.9 \\
 &  & $f_0$ $\times\, c_t(p_t^{[0]})$ & -12.3 & -8.5 & -5.6 & \textbf{-1.1} & 1.1 & 3.5 & 2.9 & -2.9 \\
 &  & $f_0$ $\times\, c_t(q_t^{[0,0]})$ & 12.4 & 8.4 & -2.3 & -18.2 & -11.7 & -8.3 & -24.4 & -6.3 \\
 &  & $f_0$ $\times\, c_t(\rho_t^{[0]})$ & 4.0 & 7.8 & 4.3 & -11.5 & -5.7 & \textbf{-1.8} & -14.1 & -2.4 \\
 &  & $f_0$ $\times\, c_t(r_t^{[0]})$ & -16.9 & -5.5 & -2.0 & -9.8 & 1.0 & 8.0 & 11.6 & -2.0 \\
 &  & $f_0$ $\times\, c_t(b_t)$ & \textbf{0.1} & 2.0 & \textbf{1.9} & -4.7 & \textbf{-0.8} & 3.0 & \textbf{1.8} & \textbf{0.5} \\
\cline{2-11}
 & \multirow{7}{*}{C: Normal only} & $M_0$ fit on $\Delta b$ & 3.5 & 5.2 & 3.3 & 3.4 & 5.1 & 3.6 & 4.5 & 4.1 \\
 &  & $f_0$ $\times\, c_t(\mathbf{z}_t^{[0,0]})$ & 1.6 & \textbf{1.1} & 5.3 & 3.3 & 4.0 & 5.8 & 8.6 & 4.2 \\
 &  & $f_0$ $\times\, c_t(p_t^{[0]})$ & -11.3 & -6.6 & 4.4 & 3.3 & 3.3 & 4.6 & \textbf{-2.8} & \textbf{-0.7} \\
 &  & $f_0$ $\times\, c_t(q_t^{[0,0]})$ & 5.9 & -1.8 & -6.9 & -7.1 & -5.5 & -10.3 & -9.0 & -5.0 \\
 &  & $f_0$ $\times\, c_t(\rho_t^{[0]})$ & 2.9 & 1.5 & -2.3 & -4.0 & -3.3 & -6.8 & -8.7 & -3.0 \\
 &  & $f_0$ $\times\, c_t(r_t^{[0]})$ & -17.8 & -11.7 & -7.7 & -8.6 & -6.4 & \textbf{1.1} & 14.7 & -5.2 \\
 &  & $f_0$ $\times\, c_t(b_t)$ & \textbf{1.1} & 1.8 & \textbf{1.0} & \textbf{0.8} & \textbf{2.6} & 2.9 & 3.4 & 1.9 \\
\end{tabular}

%     \caption{Mean signed error of the same forecasts of $\Delta P_0$,
%     showing how much of each correction's gain is the removal of a systematic
%     over-prediction rather than a tightening of the scatter. The refitted
%     forecaster is again omitted, for the reason given in
%     \Cref{tab:forecast-bias}.}
%     \label{tab:forecast-correction-bias}
% \end{table}


% \begin{table}[htbp]
%     \centering
%     \small
%     \setlength{\tabcolsep}{4pt}
%     % Generated by method.visualization.make_plots -- do not edit.
% Medians over the distinct training runs of each arm; entry = optimiser step 1, where the learning rate is still 0.
\begin{tabular}{l|l|rrrr}
 &  & \multicolumn{4}{c}{Final fine-tuning step} \\
Final dataset $X$ & Arm & Loss (entry) & Loss (median) & $\lVert g \rVert$ (entry) & $\lVert g \rVert$ (median) \\
\hline
\multirow{3}{*}{Hallucination (I)} & Normal $\times$2 ($N\,N\,X$) & 1.61 & 1.17 & 2.04 & 1.23 \\
 & Different ($X'\,N\,X$) & 1.42 & 1.15 & 1.56 & 1.17 \\
 & Same ($X\,N\,X$) & \textbf{1.26} & \textbf{1.13} & \textbf{1.22} & \textbf{1.02} \\
\hline
\multirow{3}{*}{Opinions (Mistake I)} & Normal $\times$2 ($N\,N\,X$) & 2.38 & 1.37 & 3.02 & 1.89 \\
 & Different ($X'\,N\,X$) & 2.15 & 1.37 & 2.52 & 1.87 \\
 & Same ($X\,N\,X$) & \textbf{1.63} & \textbf{1.31} & \textbf{1.45} & \textbf{1.32} \\
\hline
\multirow{3}{*}{GSM8K (Mistake II)} & Normal $\times$2 ($N\,N\,X$) & 1.04 & 0.62 & 1.75 & 0.81 \\
 & Different ($X'\,N\,X$) & 0.99 & 0.62 & 1.30 & 0.78 \\
 & Same ($X\,N\,X$) & \textbf{0.68} & \textbf{0.59} & \textbf{0.74} & \textbf{0.58} \\
\end{tabular}

%     \caption{Where each trajectory's final fine-tuning started, and where it sat over the run as a whole; the numbers behind \Cref{fig:exp3-training-curves}. ``Entry'' is the value at the first optimiser step, logged while the learning rate is still zero, so it measures the state the run was handed rather than anything the run has yet done. ``Median'' is the same quantity over all of the run's optimiser steps. $\lVert g \rVert$ is the gradient norm. Each cell is the median over the 9--10 distinct fine-tuning events it covers; because adapters are content-addressed, a fine-tuning event shared by two trajectories is counted once. Bold marks the smallest value in each column within a dataset.}
%     \label{tab:exp3-training-curves}
% \end{table}

\chapter{Forecast Error and Regret During Each Step}
\newpage

\begin{table}[htbp]
    \centering
    \singlespacing
    \scriptsize
    \setlength{\tabcolsep}{1pt}
    % Generated by method.visualization.make_plots -- do not edit.
% RMSE of the frozen forecaster and its regret against the refit at $t$, in judge points; $f_0$ fitted on the target matched to each projection.
\begin{tabular}{l|l|l|rrrrrrr|r|rrrrrrr|r}
 &  &  & \multicolumn{8}{c|}{RMSE of $f_0$} & \multicolumn{8}{c}{Regret} \\
Trait & Trunk & Projection & 0 & 1 & 2 & 3 & 4 & 5 & 6 & Mean & 0 & 1 & 2 & 3 & 4 & 5 & 6 & Mean \\
\hline
\multirow{21}{*}{Evil} & \multirow{7}{*}{A: II drivers} & $\Delta P_0$ & 10.9 & \textbf{12.3} & 10.3 & 19.3 & \textbf{18.7} & 18.3 & 18.5 & \textbf{15.5} & 6.5 & 5.0 & 5.6 & 5.3 & 4.8 & 4.2 & 4.3 & 5.1 \\
 &  & $\Delta P_t^{0\leftarrow 0,[0]}$ & 10.9 & 13.4 & \textbf{10.2} & 20.8 & 21.0 & 21.2 & 22.0 & 17.1 & 6.5 & 4.3 & 4.8 & 4.3 & 4.1 & 4.3 & 4.2 & 4.7 \\
 &  & $\Delta P_t^{t\leftarrow 0,[0]}$ & 10.9 & 13.4 & \textbf{10.2} & 19.9 & 20.5 & 20.4 & 21.4 & 16.7 & 6.5 & 3.2 & 4.0 & 2.8 & \textbf{3.1} & \textbf{3.0} & \textbf{3.2} & \textbf{3.7} \\
 &  & $\Delta P_t^{t\leftarrow t,[0]}$ & \textbf{10.8} & 14.3 & 12.1 & 24.5 & 26.6 & 25.4 & 26.2 & 20.0 & 6.5 & \textbf{2.9} & 3.8 & 5.9 & 7.3 & 6.4 & 6.2 & 5.6 \\
 &  & $\Delta P_t^{0\leftarrow 0,[t]}$ & 10.9 & 26.9 & 11.8 & 19.1 & 19.7 & 17.2 & 16.2 & 17.4 & 6.5 & 17.7 & 4.2 & 3.5 & 6.4 & 5.1 & 4.8 & 6.9 \\
 &  & $\Delta P_t^{t\leftarrow 0,[t]}$ & 10.9 & 27.8 & 11.5 & \textbf{17.6} & 19.0 & \textbf{17.0} & \textbf{15.2} & 17.0 & 6.5 & 18.3 & 3.7 & 2.3 & 6.3 & 5.6 & 4.9 & 6.8 \\
 &  & $\Delta P_t^{t\leftarrow t,[t]}$ & \textbf{10.8} & 28.4 & 11.7 & 18.1 & 20.7 & 17.5 & 17.1 & 17.8 & 6.5 & 18.6 & \textbf{3.2} & \textbf{1.9} & 7.9 & 3.8 & 4.4 & 6.6 \\
\cline{2-19}
 & \multirow{7}{*}{B: I drivers} & $\Delta P_0$ & 10.9 & 18.8 & 17.4 & \textbf{14.1} & \textbf{14.1} & \textbf{13.6} & \textbf{9.0} & \textbf{14.0} & 6.5 & 7.4 & 7.7 & 5.8 & 6.8 & 5.8 & 6.1 & 6.6 \\
 &  & $\Delta P_t^{0\leftarrow 0,[0]}$ & 10.9 & 18.5 & 19.6 & 15.6 & 16.7 & 17.3 & 10.6 & 15.6 & 6.5 & 4.8 & 7.0 & 5.3 & 6.0 & 5.6 & 6.1 & 5.9 \\
 &  & $\Delta P_t^{t\leftarrow 0,[0]}$ & 10.9 & 19.1 & 19.7 & 15.5 & 15.9 & 16.2 & 9.2 & 15.2 & 6.5 & 4.8 & 6.7 & 5.0 & 4.9 & 4.3 & \textbf{4.5} & 5.2 \\
 &  & $\Delta P_t^{t\leftarrow t,[0]}$ & \textbf{10.8} & 31.6 & 27.3 & 23.9 & 24.7 & 24.4 & 15.4 & 22.6 & 6.5 & 13.8 & 12.0 & 10.4 & 10.7 & 10.1 & 7.5 & 10.1 \\
 &  & $\Delta P_t^{0\leftarrow 0,[t]}$ & 10.9 & 17.3 & 14.4 & 16.1 & 16.2 & 16.7 & 13.8 & 15.1 & 6.5 & 2.7 & 6.4 & 5.6 & 4.2 & 4.3 & 6.3 & 5.1 \\
 &  & $\Delta P_t^{t\leftarrow 0,[t]}$ & 10.9 & \textbf{17.1} & \textbf{13.7} & 15.8 & 15.9 & 16.4 & 12.8 & 14.7 & 6.5 & 2.2 & 6.3 & 4.9 & 3.8 & 3.8 & 5.8 & 4.8 \\
 &  & $\Delta P_t^{t\leftarrow t,[t]}$ & \textbf{10.8} & 18.8 & 15.2 & 20.3 & 17.9 & 18.3 & 14.5 & 16.5 & 6.5 & \textbf{1.7} & \textbf{4.6} & \textbf{4.6} & \textbf{2.8} & \textbf{3.0} & 4.9 & \textbf{4.0} \\
\cline{2-19}
 & \multirow{7}{*}{C: Normal only} & $\Delta P_0$ & 10.9 & 11.3 & 10.7 & 10.7 & 10.5 & 10.5 & 10.2 & 10.7 & 6.5 & 6.6 & 6.5 & 5.4 & 6.5 & 6.7 & 5.9 & 6.3 \\
 &  & $\Delta P_t^{0\leftarrow 0,[0]}$ & 10.9 & 11.0 & 10.9 & 10.7 & 10.7 & 11.1 & 10.8 & 10.9 & 6.5 & 6.7 & 6.3 & 5.3 & 6.3 & 7.4 & 6.5 & 6.4 \\
 &  & $\Delta P_t^{t\leftarrow 0,[0]}$ & 10.9 & \textbf{10.8} & \textbf{10.6} & \textbf{10.5} & \textbf{10.3} & \textbf{10.1} & \textbf{9.6} & \textbf{10.4} & 6.5 & 6.2 & 5.5 & 4.6 & 5.4 & 6.1 & 4.8 & 5.6 \\
 &  & $\Delta P_t^{t\leftarrow t,[0]}$ & \textbf{10.8} & 14.9 & 15.0 & 16.0 & 16.1 & 14.8 & 13.7 & 14.5 & 6.5 & 7.5 & 6.8 & 6.6 & 7.7 & 7.5 & 5.5 & 6.9 \\
 &  & $\Delta P_t^{0\leftarrow 0,[t]}$ & 10.9 & 12.3 & 12.8 & 12.8 & 13.2 & 12.9 & 13.0 & 12.6 & 6.5 & 5.1 & 6.8 & 6.6 & 7.7 & 6.6 & 5.5 & 6.4 \\
 &  & $\Delta P_t^{t\leftarrow 0,[t]}$ & 10.9 & 11.9 & 12.4 & 12.5 & 12.6 & 12.1 & 12.3 & 12.1 & 6.5 & 4.6 & 5.9 & 6.0 & 6.9 & 6.0 & 4.6 & 5.8 \\
 &  & $\Delta P_t^{t\leftarrow t,[t]}$ & \textbf{10.8} & 12.2 & 13.5 & 14.3 & 14.5 & 12.7 & 12.6 & 13.0 & 6.5 & \textbf{3.7} & \textbf{3.9} & \textbf{3.7} & \textbf{4.5} & \textbf{4.4} & \textbf{2.9} & \textbf{4.2} \\
\hline
\multirow{21}{*}{Sycophancy} & \multirow{7}{*}{A: II drivers} & $\Delta P_0$ & 15.2 & 26.9 & 15.2 & 23.5 & 20.8 & 42.5 & 40.0 & 26.3 & 3.0 & 7.0 & 3.1 & 6.4 & 4.9 & 21.3 & 19.6 & 9.3 \\
 &  & $\Delta P_t^{0\leftarrow 0,[0]}$ & 15.2 & 25.0 & \textbf{12.1} & 21.5 & 19.8 & 42.2 & 41.1 & 25.3 & 3.0 & 5.4 & \textbf{2.0} & 5.8 & 5.3 & 20.6 & 20.4 & 8.9 \\
 &  & $\Delta P_t^{t\leftarrow 0,[0]}$ & 15.2 & 27.3 & 13.0 & 24.1 & 22.4 & 42.9 & 42.3 & 26.7 & 3.0 & 7.0 & 2.8 & 6.9 & 6.8 & 21.2 & 21.4 & 9.9 \\
 &  & $\Delta P_t^{t\leftarrow t,[0]}$ & \textbf{14.9} & 28.8 & 13.8 & 28.0 & 27.8 & 46.3 & 46.0 & 29.4 & \textbf{2.8} & 8.2 & 2.4 & 10.1 & 10.7 & 24.8 & 25.3 & 12.1 \\
 &  & $\Delta P_t^{0\leftarrow 0,[t]}$ & 15.2 & 20.8 & 13.0 & 13.4 & \textbf{15.2} & 16.9 & \textbf{28.3} & 17.5 & 3.0 & 4.8 & 3.7 & 2.6 & \textbf{0.7} & 4.8 & 8.1 & 4.0 \\
 &  & $\Delta P_t^{t\leftarrow 0,[t]}$ & 15.2 & \textbf{19.3} & 13.0 & \textbf{10.6} & 16.4 & \textbf{14.5} & 28.8 & \textbf{16.8} & 3.0 & \textbf{3.5} & 4.1 & \textbf{0.1} & 2.1 & \textbf{1.1} & 8.3 & \textbf{3.2} \\
 &  & $\Delta P_t^{t\leftarrow t,[t]}$ & \textbf{14.9} & 22.3 & 14.2 & 19.1 & 20.6 & 18.7 & 28.8 & 19.8 & \textbf{2.8} & \textbf{3.5} & 3.4 & 2.2 & 2.9 & 1.6 & \textbf{7.8} & 3.5 \\
\cline{2-19}
 & \multirow{7}{*}{B: I drivers} & $\Delta P_0$ & 15.2 & 16.5 & 19.8 & 13.9 & 16.3 & 18.2 & 13.7 & 16.2 & 3.0 & 3.0 & 3.7 & 2.0 & 3.4 & 3.2 & 3.6 & 3.1 \\
 &  & $\Delta P_t^{0\leftarrow 0,[0]}$ & 15.2 & \textbf{12.2} & 15.4 & \textbf{10.9} & 13.4 & 15.5 & 9.6 & 13.2 & 3.0 & \textbf{1.3} & 2.2 & \textbf{1.9} & 3.0 & 2.9 & 1.9 & \textbf{2.3} \\
 &  & $\Delta P_t^{t\leftarrow 0,[0]}$ & 15.2 & 14.4 & 17.2 & 12.3 & 15.2 & 16.7 & \textbf{9.0} & 14.3 & 3.0 & 2.6 & 3.6 & 2.6 & 4.1 & 3.7 & \textbf{1.3} & 3.0 \\
 &  & $\Delta P_t^{t\leftarrow t,[0]}$ & \textbf{14.9} & 19.0 & 20.4 & 22.6 & 22.4 & 23.0 & 19.9 & 20.3 & \textbf{2.8} & 6.6 & 6.8 & 10.3 & 10.4 & 9.6 & 6.3 & 7.5 \\
 &  & $\Delta P_t^{0\leftarrow 0,[t]}$ & 15.2 & 13.4 & \textbf{14.7} & 13.9 & 10.8 & 10.4 & 9.4 & 12.5 & 3.0 & 4.8 & \textbf{1.7} & 4.4 & \textbf{2.6} & \textbf{0.3} & 1.8 & 2.7 \\
 &  & $\Delta P_t^{t\leftarrow 0,[t]}$ & 15.2 & 12.8 & 14.8 & 14.3 & \textbf{10.2} & \textbf{9.1} & 9.4 & \textbf{12.3} & 3.0 & 4.4 & 1.8 & 4.4 & 2.8 & 0.4 & 2.4 & 2.8 \\
 &  & $\Delta P_t^{t\leftarrow t,[t]}$ & \textbf{14.9} & 14.0 & 15.7 & 19.3 & 13.5 & 12.0 & 11.6 & 14.4 & \textbf{2.8} & 4.0 & 4.6 & 5.4 & 3.5 & 0.6 & 1.4 & 3.2 \\
\cline{2-19}
 & \multirow{7}{*}{C: Normal only} & $\Delta P_0$ & 15.2 & 13.2 & 13.6 & 12.7 & 14.1 & 14.0 & 17.3 & 14.3 & 3.0 & 2.6 & 2.9 & 2.6 & 3.5 & 3.7 & 3.5 & 3.1 \\
 &  & $\Delta P_t^{0\leftarrow 0,[0]}$ & 15.2 & 10.5 & \textbf{10.7} & \textbf{9.9} & \textbf{10.3} & 10.4 & 14.7 & \textbf{11.7} & 3.0 & \textbf{1.7} & \textbf{1.8} & \textbf{2.1} & \textbf{1.6} & 2.4 & 2.0 & \textbf{2.1} \\
 &  & $\Delta P_t^{t\leftarrow 0,[0]}$ & 15.2 & \textbf{10.4} & 11.0 & 10.3 & 10.7 & 9.7 & 15.1 & 11.8 & 3.0 & 1.8 & 2.1 & 2.6 & 2.0 & 1.7 & \textbf{1.7} & \textbf{2.1} \\
 &  & $\Delta P_t^{t\leftarrow t,[0]}$ & \textbf{14.9} & 16.2 & 18.9 & 22.1 & 20.5 & 17.9 & 18.1 & 18.4 & \textbf{2.8} & 6.8 & 8.1 & 10.5 & 9.1 & 5.8 & 2.2 & 6.5 \\
 &  & $\Delta P_t^{0\leftarrow 0,[t]}$ & 15.2 & 12.3 & 12.5 & 13.4 & 13.2 & 9.8 & 12.1 & 12.6 & 3.0 & 3.7 & 2.7 & 3.8 & 2.9 & 1.8 & 1.8 & 2.8 \\
 &  & $\Delta P_t^{t\leftarrow 0,[t]}$ & 15.2 & 12.8 & 12.8 & 13.6 & 13.2 & \textbf{9.4} & \textbf{11.5} & 12.6 & 3.0 & 4.2 & 3.2 & 4.3 & 3.3 & 2.1 & 2.3 & 3.2 \\
 &  & $\Delta P_t^{t\leftarrow t,[t]}$ & \textbf{14.9} & 13.6 & 14.8 & 15.6 & 16.1 & 10.7 & 16.9 & 14.6 & \textbf{2.8} & 3.0 & 2.0 & \textbf{2.1} & 2.1 & \textbf{0.8} & 3.4 & 2.3 \\
\end{tabular}

    \caption{Out-of-sample error of the frozen forecaster $f_0$, and its
    regret, in judge points, at every checkpoint of every trunk for each of
    the projection difference variants discussed in
    \Cref{sec:deltap-variants}. Regret is the gap to the forecaster refitted
    at that same checkpoint, $\text{RMSE}(f_0) - \text{RMSE}(f_t)$. Both quantities
    use the leave-one-out protocol of \Cref{subsec:forecaster_eval}, and $f_0$
    is fitted on the target matched to each projection, as in
    \Cref{fig:exp2-headline-rmse}: $\Delta b_{t+1}$ for the variants that
    regenerate the answers behind $\mathbf{h}^{\text{predicted}}_t$ and
    $b_{t+1}$ for those that keep $\mathcal{M}_0$'s.}
    \label{tab:forecast-regret}
\end{table}
